\section{Background and General Goals}
\label{sec:background}

\subsection{Background}

Data extraction means giving an AI a document, like an invoice or a registration form, and asking it to find and copy out specific pieces of information. For example: ``What is the total amount?'' or ``What is the applicant's name?'' This kind of task is needed every day in schools, hospitals, banks, and government offices where large numbers of forms and documents must be processed quickly and correctly.

Right now, most organizations use a paid online AI service to do this. They upload their documents to the internet, the AI reads them, and sends back the answers. This works well, but every upload costs a small amount of money. When an organization processes thousands of documents every day, those small costs become a very large monthly expense. On top of that, sending private documents to someone else's server over the internet creates a serious privacy risk, especially in sensitive fields like healthcare or banking.

A newer option is now available. Researchers have found ways to make large AI models much smaller so they can run on a regular laptop without needing the internet at all \cite{dettmers2022llmint8, frantar2022gptq}. A free tool called Ollama \cite{ollama_2024} makes this even easier by handling all the downloading and setup automatically --- you just type the model name and Ollama does the rest. Three local models will be used in this study: \texttt{deepseek-ocr} by DeepSeek AI \cite{deepseek_ocr_2024}, which focuses on reading text from document images; \texttt{gemma4:e4b} by Google DeepMind \cite{gemma4_2026}, which can process document images without a separate reading tool; and \texttt{qwen3.5:4b} by Alibaba \cite{qwen3_2025}, which is strong at pulling out structured information like form fields. A fourth model, \texttt{Phi-4-multimodal-instruct} by Microsoft \cite{phi4mini_2025}, will also be tested as an optional extra if Ollama adds support for it. Each model is built in a different way, which is why testing all three is important --- it means the results will apply to local AI tools in general, not just one specific model \cite{gemma4_phi4_qwen3_2026}.

The problem is that nobody has done a proper, honest test to find out whether any of these local models are good enough to replace a paid online service --- or whether the results are affected by how the instruction is worded.

\subsection{Statement of Objectives}

This thesis will do three things:

\begin{enumerate}
    \item Test and compare how accurately three local AI models (\texttt{deepseek-ocr}, \texttt{gemma4:e4b}, and \texttt{qwen3.5:4b}, all run via Ollama) and two paid online AI services (GPT-4o by OpenAI \cite{openai_gpt4o_2024} and Gemini 3.1 Pro by Google \cite{google_gemini_2025}) extract information from the same set of documents. Each model will be tested using three different ways of writing the instruction, so that the results are fair and do not depend on how the question happened to be worded \cite{anlsstar_2024, gemma4_phi4_qwen3_2026}.
    \item Measure and compare how fast each option is --- specifically how long it takes from sending a document to receiving the extracted information.
    \item Calculate the exact cost per document for each option. For local models, the electricity used during each test will be measured with a software tool (\texttt{powermetrics} on Mac or \texttt{powertop} on Linux), and the laptop's purchase price will be spread over three years using straight-line depreciation \cite{lenovo_tco_2026}. For the online service, the cost will come from the billing records. These numbers will then be used to find the number of daily documents at which a local model becomes cheaper than the online service.
\end{enumerate}

\subsection{Significance of the Project}

This research answers a question that real organizations are facing right now. As privacy laws get stricter and online AI costs keep growing, choosing between a local AI tool and a paid online service has very real financial and legal consequences. This thesis will give organizations clear, honest numbers rather than marketing claims, so they can make a smarter decision that fits their own situation. The findings will also add new knowledge to the field of AI deployment, which is an area that is growing fast but still lacks practical, independent research.
